% =====================================================================================================================
% Appendix results: % =====================================================================================================================
% Appendix results: % =====================================================================================================================
% Appendix results: % =====================================================================================================================
% Appendix results: \input{appendix_results} in the appendix. Files and numbers as in experiments_results.tex.
% Order is by importance; drop subsections from the end if the appendix must be short.
% =====================================================================================================================
\providecommand{\figdir}{../paper_graphs}

\section{Additional Results}\label{app:results}

% ---------------------------------------------------------------------------------------------------------------------
\subsection{Per-model results}\label{app:per_model}

\begin{table}[htbp]
  \centering
  \caption{Table~\ref{tab:pooled} on the $c=1.2$ world models.}
  \label{tab:pooled_c1_2}
  \small
  \input{\figdir/appendix/table_pooled_c1_2}
\end{table}

\begin{table}[htbp]
  \centering
  \caption{Table~\ref{tab:pooled} on the unbounded ($c=\infty$) world models.}
  \label{tab:pooled_cinf}
  \small
  \input{\figdir/appendix/table_pooled_cinf}
\end{table}

Table~\ref{tab:per_model} gives each planner's violation rate and return on every $N=10$ world model, pooled over
$\epsilon\ge0.1$, and Figs.~\ref{fig:grid_violations} and~\ref{fig:grid_return} show the full sweeps. Pure SCP is
significantly less safe than SCP+MPPI on every model. Against the samplers, the only significant difference is MPPI
with 512 sequences on PointGoal2 with $c=1.2$, which violates less.

\begin{table}[htbp]
  \centering
  \caption{Each planner on each $N=10$ world model, pooled over $\epsilon\ge0.1$: violation rate (Wilson) and return
  (bootstrap) with 95\% intervals, and McNemar's $p$ for SCP+MPPI against it on paired seeds, Holm-corrected over the
  six models.}
  \label{tab:per_model}
  \small
  \input{\figdir/appendix/table_per_model}
\end{table}

\begin{figure}[htbp]
  \centering
  \includegraphics[width=\linewidth]{\figdir/appendix/fig_eps_sweep_all_bounds_violations.pdf}
  \caption{Violation rate against $\epsilon$ for every planner on every $N=10$ world model (rows: tasks; columns:
  Lipschitz bound $c$), with 95\% Wilson intervals.}
  \label{fig:grid_violations}
\end{figure}

\begin{figure}[htbp]
  \centering
  \includegraphics[width=\linewidth]{\figdir/appendix/fig_eps_sweep_all_bounds_return.pdf}
  \caption{Return against $\epsilon$ for every planner on every $N=10$ world model, with 95\% bootstrap intervals.}
  \label{fig:grid_return}
\end{figure}

% ---------------------------------------------------------------------------------------------------------------------
\subsection{Lipschitz bound and SCP's convergence}\label{app:lipschitz}

\begin{figure}[htbp]
  \centering
  \includegraphics[width=\linewidth]{\figdir/appendix/fig_lipschitz_convergence.pdf}
  \caption{Share of SCP+MPPI's steps that reached the 40-iteration cap against $\epsilon$, for each Lipschitz bound,
  at $N=10$.}
  \label{fig:lipschitz_convergence}
\end{figure}

Without a Lipschitz bound, SCP+MPPI reaches the iteration cap on more steps than with either bound at every margin,
except at $\epsilon=0.5$ on PointGoal2, where $c=1.5$ and $c=\infty$ are within 0.02 (Fig.~\ref{fig:lipschitz_convergence}).
Which bound converges best differs by task: $c=1.5$ on PointGoal1, $c=1.2$ on PointGoal2.

% ---------------------------------------------------------------------------------------------------------------------
\subsection{What more samples buy at \texorpdfstring{$N=10$}{N=10}}\label{app:samples}

\begin{figure}[htbp]
  \centering
  \includegraphics[width=\linewidth]{\figdir/ablation/fig_sample_vs_time_matched.pdf}
  \caption{Fig.~\ref{fig:eps_sweep} with the samplers at both budgets: 512 sequences (open markers, dashed) and
  time-matched (filled, solid).}
  \label{fig:sample_vs_time_matched}
\end{figure}

\begin{figure}[htbp]
  \centering
  \includegraphics[width=\linewidth]{\figdir/ablation/fig_return_vs_violations.pdf}
  \caption{Return against violation rate on each task's $c=1.5$ model at $N=10$, pooled over $\epsilon\ge0.1$, with
  95\% intervals. Arrows lead from each sampler with 512 sequences to the same sampler time-matched.}
  \label{fig:trade_off}
\end{figure}

Figs.~\ref{fig:sample_vs_time_matched} and~\ref{fig:trade_off} show what 128 times as many sequences buy the samplers.
On PointGoal1's $c=1.5$ model, the time-matched samplers earn more than with 512 sequences but violate more often
(MPPI: return 0.78 to 1.15, violation rate 0.19 to 0.24; CEM: 0.84 to 0.94 and 0.14 to 0.26), while SCP+MPPI reaches
the higher return at the lower violation rate (1.15 and 0.17). On PointGoal2 the extra sequences raise return (0.35 to
0.51 and 0.49) without changing the violation rate. These shifts lie within the intervals; the paired tests of
Sec.~\ref{sec:q1} find no significant difference in safety between SCP+MPPI and either budget.

% ---------------------------------------------------------------------------------------------------------------------
\subsection{All budgets at every horizon}\label{app:horizon_budgets}

\begin{figure}[htbp]
  \centering
  \includegraphics[width=\linewidth]{\figdir/ablation/fig_horizon_all_budgets.pdf}
  \caption{Fig.~\ref{fig:horizon} with the samplers at both 512 sequences (open markers, dashed) and their
  time-matched budget (filled, solid).}
  \label{fig:horizon_all}
\end{figure}

With 512 sequences the samplers' return at $N=20$ is close to SCP+MPPI's (MPPI 2.07, CEM 2.58, against 2.34; neither
difference significant), and at $N=30$ below it (1.59 and 1.22 against 2.26; $p=0.043$ and $0.013$)
(Fig.~\ref{fig:horizon_all}). Together with Fig.~\ref{fig:compute}, this shows that the samplers' best budget depends on
the sampler and the horizon.

% ---------------------------------------------------------------------------------------------------------------------
\subsection{Does optimisation exploit the model?}\label{app:optimism}

\begin{figure}[htbp]
  \centering
  \includegraphics[width=\linewidth]{\figdir/ablation/fig_model_optimism.pdf}
  \caption{Model optimism: each executed plan's predicted return (discounted predicted rewards plus terminal value)
  minus the discounted return received from that step on, averaged over each episode, with 95\% bootstrap intervals.
  $N=10$ pools every world model ($\epsilon\ge0.1$); $N=20$ and $30$ are the horizon study's models.}
  \label{fig:optimism}
\end{figure}

A gradient-based optimiser might exploit model error, finding plans whose predicted return the environment does not
deliver. We measure each executed plan's optimism: its predicted return, whose terminal value carries the goal bonus,
minus the discounted return actually received from that step on, using steps at least 299 steps before the episode
ends so that the part of the future the episode did not reach weighs under 5\%. Every planner is optimistic by about
2, as expected of a value fitted to the 0.7 expectile, but by nearly the same amount (Fig.~\ref{fig:optimism}): 2.06
(pure SCP) to 2.19 at $N=10$, 2.20 to 2.44 at $N=20$ and 2.41 to 2.56 at $N=30$, with SCP+MPPI inside the samplers'
range at every horizon and no trend with the sampling budget. We find no evidence that SCP, or a larger sampling
budget, exploits the model more than sampling with 512 sequences does.

% ---------------------------------------------------------------------------------------------------------------------
\subsection{Warned and unwarned violations}\label{app:crashes}

\begin{figure}[htbp]
  \centering
  \includegraphics[width=\linewidth]{\figdir/ablation/fig_crash_types.pdf}
  \caption{Violations per 100 episodes over the six $N=10$ models and $\epsilon\ge0.1$, split into those preceded
  within $N$ steps by a plan that broke the barrier in the model (warned, light) and those that were not (unwarned,
  dark, counts above the bars).}
  \label{fig:crash_types}
\end{figure}

Few violations come without warning (Fig.~\ref{fig:crash_types}). SCP+MPPI has 14 unwarned violations in 880
episodes, against 6 and 7 for the 512-sequence samplers and 7 and 13 for the time-matched ones. That stronger
optimisation, by SCP or by more samples, leads to more violations the model did not foresee would be consistent with
it planning closer to where the model is wrong, but paired by seed none of these differences is significant
($p\ge0.12$).

% ---------------------------------------------------------------------------------------------------------------------
\subsection{When violations happen}\label{app:survival}

\begin{figure}[htbp]
  \centering
  \includegraphics[width=\linewidth]{\figdir/appendix/fig_survival_pg1_c1_5.pdf}
  \caption{Share of episodes without a violation by each step, at each margin, on PointGoal1's $c=1.5$ model at
  $N=10$. Episodes end early only at a violation, so these are Kaplan--Meier curves without censoring.}
  \label{fig:survival}
\end{figure}

Most violations happen within the first 300 or so steps, after which the curves flatten (Fig.~\ref{fig:survival});
pure SCP's keeps falling for longer.
 in the appendix. Files and numbers as in experiments_results.tex.
% Order is by importance; drop subsections from the end if the appendix must be short.
% =====================================================================================================================
\providecommand{\figdir}{../paper_graphs}

\section{Additional Results}\label{app:results}

% ---------------------------------------------------------------------------------------------------------------------
\subsection{Per-model results}\label{app:per_model}

\begin{table}[htbp]
  \centering
  \caption{Table~\ref{tab:pooled} on the $c=1.2$ world models.}
  \label{tab:pooled_c1_2}
  \small
  \input{\figdir/appendix/table_pooled_c1_2}
\end{table}

\begin{table}[htbp]
  \centering
  \caption{Table~\ref{tab:pooled} on the unbounded ($c=\infty$) world models.}
  \label{tab:pooled_cinf}
  \small
  \input{\figdir/appendix/table_pooled_cinf}
\end{table}

Table~\ref{tab:per_model} gives each planner's violation rate and return on every $N=10$ world model, pooled over
$\epsilon\ge0.1$, and Figs.~\ref{fig:grid_violations} and~\ref{fig:grid_return} show the full sweeps. Pure SCP is
significantly less safe than SCP+MPPI on every model. Against the samplers, the only significant difference is MPPI
with 512 sequences on PointGoal2 with $c=1.2$, which violates less.

\begin{table}[htbp]
  \centering
  \caption{Each planner on each $N=10$ world model, pooled over $\epsilon\ge0.1$: violation rate (Wilson) and return
  (bootstrap) with 95\% intervals, and McNemar's $p$ for SCP+MPPI against it on paired seeds, Holm-corrected over the
  six models.}
  \label{tab:per_model}
  \small
  \input{\figdir/appendix/table_per_model}
\end{table}

\begin{figure}[htbp]
  \centering
  \includegraphics[width=\linewidth]{\figdir/appendix/fig_eps_sweep_all_bounds_violations.pdf}
  \caption{Violation rate against $\epsilon$ for every planner on every $N=10$ world model (rows: tasks; columns:
  Lipschitz bound $c$), with 95\% Wilson intervals.}
  \label{fig:grid_violations}
\end{figure}

\begin{figure}[htbp]
  \centering
  \includegraphics[width=\linewidth]{\figdir/appendix/fig_eps_sweep_all_bounds_return.pdf}
  \caption{Return against $\epsilon$ for every planner on every $N=10$ world model, with 95\% bootstrap intervals.}
  \label{fig:grid_return}
\end{figure}

% ---------------------------------------------------------------------------------------------------------------------
\subsection{Lipschitz bound and SCP's convergence}\label{app:lipschitz}

\begin{figure}[htbp]
  \centering
  \includegraphics[width=\linewidth]{\figdir/appendix/fig_lipschitz_convergence.pdf}
  \caption{Share of SCP+MPPI's steps that reached the 40-iteration cap against $\epsilon$, for each Lipschitz bound,
  at $N=10$.}
  \label{fig:lipschitz_convergence}
\end{figure}

Without a Lipschitz bound, SCP+MPPI reaches the iteration cap on more steps than with either bound at every margin,
except at $\epsilon=0.5$ on PointGoal2, where $c=1.5$ and $c=\infty$ are within 0.02 (Fig.~\ref{fig:lipschitz_convergence}).
Which bound converges best differs by task: $c=1.5$ on PointGoal1, $c=1.2$ on PointGoal2.

% ---------------------------------------------------------------------------------------------------------------------
\subsection{What more samples buy at \texorpdfstring{$N=10$}{N=10}}\label{app:samples}

\begin{figure}[htbp]
  \centering
  \includegraphics[width=\linewidth]{\figdir/ablation/fig_sample_vs_time_matched.pdf}
  \caption{Fig.~\ref{fig:eps_sweep} with the samplers at both budgets: 512 sequences (open markers, dashed) and
  time-matched (filled, solid).}
  \label{fig:sample_vs_time_matched}
\end{figure}

\begin{figure}[htbp]
  \centering
  \includegraphics[width=\linewidth]{\figdir/ablation/fig_return_vs_violations.pdf}
  \caption{Return against violation rate on each task's $c=1.5$ model at $N=10$, pooled over $\epsilon\ge0.1$, with
  95\% intervals. Arrows lead from each sampler with 512 sequences to the same sampler time-matched.}
  \label{fig:trade_off}
\end{figure}

Figs.~\ref{fig:sample_vs_time_matched} and~\ref{fig:trade_off} show what 128 times as many sequences buy the samplers.
On PointGoal1's $c=1.5$ model, the time-matched samplers earn more than with 512 sequences but violate more often
(MPPI: return 0.78 to 1.15, violation rate 0.19 to 0.24; CEM: 0.84 to 0.94 and 0.14 to 0.26), while SCP+MPPI reaches
the higher return at the lower violation rate (1.15 and 0.17). On PointGoal2 the extra sequences raise return (0.35 to
0.51 and 0.49) without changing the violation rate. These shifts lie within the intervals; the paired tests of
Sec.~\ref{sec:q1} find no significant difference in safety between SCP+MPPI and either budget.

% ---------------------------------------------------------------------------------------------------------------------
\subsection{All budgets at every horizon}\label{app:horizon_budgets}

\begin{figure}[htbp]
  \centering
  \includegraphics[width=\linewidth]{\figdir/ablation/fig_horizon_all_budgets.pdf}
  \caption{Fig.~\ref{fig:horizon} with the samplers at both 512 sequences (open markers, dashed) and their
  time-matched budget (filled, solid).}
  \label{fig:horizon_all}
\end{figure}

With 512 sequences the samplers' return at $N=20$ is close to SCP+MPPI's (MPPI 2.07, CEM 2.58, against 2.34; neither
difference significant), and at $N=30$ below it (1.59 and 1.22 against 2.26; $p=0.043$ and $0.013$)
(Fig.~\ref{fig:horizon_all}). Together with Fig.~\ref{fig:compute}, this shows that the samplers' best budget depends on
the sampler and the horizon.

% ---------------------------------------------------------------------------------------------------------------------
\subsection{Does optimisation exploit the model?}\label{app:optimism}

\begin{figure}[htbp]
  \centering
  \includegraphics[width=\linewidth]{\figdir/ablation/fig_model_optimism.pdf}
  \caption{Model optimism: each executed plan's predicted return (discounted predicted rewards plus terminal value)
  minus the discounted return received from that step on, averaged over each episode, with 95\% bootstrap intervals.
  $N=10$ pools every world model ($\epsilon\ge0.1$); $N=20$ and $30$ are the horizon study's models.}
  \label{fig:optimism}
\end{figure}

A gradient-based optimiser might exploit model error, finding plans whose predicted return the environment does not
deliver. We measure each executed plan's optimism: its predicted return, whose terminal value carries the goal bonus,
minus the discounted return actually received from that step on, using steps at least 299 steps before the episode
ends so that the part of the future the episode did not reach weighs under 5\%. Every planner is optimistic by about
2, as expected of a value fitted to the 0.7 expectile, but by nearly the same amount (Fig.~\ref{fig:optimism}): 2.06
(pure SCP) to 2.19 at $N=10$, 2.20 to 2.44 at $N=20$ and 2.41 to 2.56 at $N=30$, with SCP+MPPI inside the samplers'
range at every horizon and no trend with the sampling budget. We find no evidence that SCP, or a larger sampling
budget, exploits the model more than sampling with 512 sequences does.

% ---------------------------------------------------------------------------------------------------------------------
\subsection{Warned and unwarned violations}\label{app:crashes}

\begin{figure}[htbp]
  \centering
  \includegraphics[width=\linewidth]{\figdir/ablation/fig_crash_types.pdf}
  \caption{Violations per 100 episodes over the six $N=10$ models and $\epsilon\ge0.1$, split into those preceded
  within $N$ steps by a plan that broke the barrier in the model (warned, light) and those that were not (unwarned,
  dark, counts above the bars).}
  \label{fig:crash_types}
\end{figure}

Few violations come without warning (Fig.~\ref{fig:crash_types}). SCP+MPPI has 14 unwarned violations in 880
episodes, against 6 and 7 for the 512-sequence samplers and 7 and 13 for the time-matched ones. That stronger
optimisation, by SCP or by more samples, leads to more violations the model did not foresee would be consistent with
it planning closer to where the model is wrong, but paired by seed none of these differences is significant
($p\ge0.12$).

% ---------------------------------------------------------------------------------------------------------------------
\subsection{When violations happen}\label{app:survival}

\begin{figure}[htbp]
  \centering
  \includegraphics[width=\linewidth]{\figdir/appendix/fig_survival_pg1_c1_5.pdf}
  \caption{Share of episodes without a violation by each step, at each margin, on PointGoal1's $c=1.5$ model at
  $N=10$. Episodes end early only at a violation, so these are Kaplan--Meier curves without censoring.}
  \label{fig:survival}
\end{figure}

Most violations happen within the first 300 or so steps, after which the curves flatten (Fig.~\ref{fig:survival});
pure SCP's keeps falling for longer.
 in the appendix. Files and numbers as in experiments_results.tex.
% Order is by importance; drop subsections from the end if the appendix must be short.
% =====================================================================================================================
\providecommand{\figdir}{../paper_graphs}

\section{Additional Results}\label{app:results}

% ---------------------------------------------------------------------------------------------------------------------
\subsection{Per-model results}\label{app:per_model}

\begin{table}[htbp]
  \centering
  \caption{Table~\ref{tab:pooled} on the $c=1.2$ world models.}
  \label{tab:pooled_c1_2}
  \small
  \input{\figdir/appendix/table_pooled_c1_2}
\end{table}

\begin{table}[htbp]
  \centering
  \caption{Table~\ref{tab:pooled} on the unbounded ($c=\infty$) world models.}
  \label{tab:pooled_cinf}
  \small
  \input{\figdir/appendix/table_pooled_cinf}
\end{table}

Table~\ref{tab:per_model} gives each planner's violation rate and return on every $N=10$ world model, pooled over
$\epsilon\ge0.1$, and Figs.~\ref{fig:grid_violations} and~\ref{fig:grid_return} show the full sweeps. Pure SCP is
significantly less safe than SCP+MPPI on every model. Against the samplers, the only significant difference is MPPI
with 512 sequences on PointGoal2 with $c=1.2$, which violates less.

\begin{table}[htbp]
  \centering
  \caption{Each planner on each $N=10$ world model, pooled over $\epsilon\ge0.1$: violation rate (Wilson) and return
  (bootstrap) with 95\% intervals, and McNemar's $p$ for SCP+MPPI against it on paired seeds, Holm-corrected over the
  six models.}
  \label{tab:per_model}
  \small
  \input{\figdir/appendix/table_per_model}
\end{table}

\begin{figure}[htbp]
  \centering
  \includegraphics[width=\linewidth]{\figdir/appendix/fig_eps_sweep_all_bounds_violations.pdf}
  \caption{Violation rate against $\epsilon$ for every planner on every $N=10$ world model (rows: tasks; columns:
  Lipschitz bound $c$), with 95\% Wilson intervals.}
  \label{fig:grid_violations}
\end{figure}

\begin{figure}[htbp]
  \centering
  \includegraphics[width=\linewidth]{\figdir/appendix/fig_eps_sweep_all_bounds_return.pdf}
  \caption{Return against $\epsilon$ for every planner on every $N=10$ world model, with 95\% bootstrap intervals.}
  \label{fig:grid_return}
\end{figure}

% ---------------------------------------------------------------------------------------------------------------------
\subsection{Lipschitz bound and SCP's convergence}\label{app:lipschitz}

\begin{figure}[htbp]
  \centering
  \includegraphics[width=\linewidth]{\figdir/appendix/fig_lipschitz_convergence.pdf}
  \caption{Share of SCP+MPPI's steps that reached the 40-iteration cap against $\epsilon$, for each Lipschitz bound,
  at $N=10$.}
  \label{fig:lipschitz_convergence}
\end{figure}

Without a Lipschitz bound, SCP+MPPI reaches the iteration cap on more steps than with either bound at every margin,
except at $\epsilon=0.5$ on PointGoal2, where $c=1.5$ and $c=\infty$ are within 0.02 (Fig.~\ref{fig:lipschitz_convergence}).
Which bound converges best differs by task: $c=1.5$ on PointGoal1, $c=1.2$ on PointGoal2.

% ---------------------------------------------------------------------------------------------------------------------
\subsection{What more samples buy at \texorpdfstring{$N=10$}{N=10}}\label{app:samples}

\begin{figure}[htbp]
  \centering
  \includegraphics[width=\linewidth]{\figdir/ablation/fig_sample_vs_time_matched.pdf}
  \caption{Fig.~\ref{fig:eps_sweep} with the samplers at both budgets: 512 sequences (open markers, dashed) and
  time-matched (filled, solid).}
  \label{fig:sample_vs_time_matched}
\end{figure}

\begin{figure}[htbp]
  \centering
  \includegraphics[width=\linewidth]{\figdir/ablation/fig_return_vs_violations.pdf}
  \caption{Return against violation rate on each task's $c=1.5$ model at $N=10$, pooled over $\epsilon\ge0.1$, with
  95\% intervals. Arrows lead from each sampler with 512 sequences to the same sampler time-matched.}
  \label{fig:trade_off}
\end{figure}

Figs.~\ref{fig:sample_vs_time_matched} and~\ref{fig:trade_off} show what 128 times as many sequences buy the samplers.
On PointGoal1's $c=1.5$ model, the time-matched samplers earn more than with 512 sequences but violate more often
(MPPI: return 0.78 to 1.15, violation rate 0.19 to 0.24; CEM: 0.84 to 0.94 and 0.14 to 0.26), while SCP+MPPI reaches
the higher return at the lower violation rate (1.15 and 0.17). On PointGoal2 the extra sequences raise return (0.35 to
0.51 and 0.49) without changing the violation rate. These shifts lie within the intervals; the paired tests of
Sec.~\ref{sec:q1} find no significant difference in safety between SCP+MPPI and either budget.

% ---------------------------------------------------------------------------------------------------------------------
\subsection{All budgets at every horizon}\label{app:horizon_budgets}

\begin{figure}[htbp]
  \centering
  \includegraphics[width=\linewidth]{\figdir/ablation/fig_horizon_all_budgets.pdf}
  \caption{Fig.~\ref{fig:horizon} with the samplers at both 512 sequences (open markers, dashed) and their
  time-matched budget (filled, solid).}
  \label{fig:horizon_all}
\end{figure}

With 512 sequences the samplers' return at $N=20$ is close to SCP+MPPI's (MPPI 2.07, CEM 2.58, against 2.34; neither
difference significant), and at $N=30$ below it (1.59 and 1.22 against 2.26; $p=0.043$ and $0.013$)
(Fig.~\ref{fig:horizon_all}). Together with Fig.~\ref{fig:compute}, this shows that the samplers' best budget depends on
the sampler and the horizon.

% ---------------------------------------------------------------------------------------------------------------------
\subsection{Does optimisation exploit the model?}\label{app:optimism}

\begin{figure}[htbp]
  \centering
  \includegraphics[width=\linewidth]{\figdir/ablation/fig_model_optimism.pdf}
  \caption{Model optimism: each executed plan's predicted return (discounted predicted rewards plus terminal value)
  minus the discounted return received from that step on, averaged over each episode, with 95\% bootstrap intervals.
  $N=10$ pools every world model ($\epsilon\ge0.1$); $N=20$ and $30$ are the horizon study's models.}
  \label{fig:optimism}
\end{figure}

A gradient-based optimiser might exploit model error, finding plans whose predicted return the environment does not
deliver. We measure each executed plan's optimism: its predicted return, whose terminal value carries the goal bonus,
minus the discounted return actually received from that step on, using steps at least 299 steps before the episode
ends so that the part of the future the episode did not reach weighs under 5\%. Every planner is optimistic by about
2, as expected of a value fitted to the 0.7 expectile, but by nearly the same amount (Fig.~\ref{fig:optimism}): 2.06
(pure SCP) to 2.19 at $N=10$, 2.20 to 2.44 at $N=20$ and 2.41 to 2.56 at $N=30$, with SCP+MPPI inside the samplers'
range at every horizon and no trend with the sampling budget. We find no evidence that SCP, or a larger sampling
budget, exploits the model more than sampling with 512 sequences does.

% ---------------------------------------------------------------------------------------------------------------------
\subsection{Warned and unwarned violations}\label{app:crashes}

\begin{figure}[htbp]
  \centering
  \includegraphics[width=\linewidth]{\figdir/ablation/fig_crash_types.pdf}
  \caption{Violations per 100 episodes over the six $N=10$ models and $\epsilon\ge0.1$, split into those preceded
  within $N$ steps by a plan that broke the barrier in the model (warned, light) and those that were not (unwarned,
  dark, counts above the bars).}
  \label{fig:crash_types}
\end{figure}

Few violations come without warning (Fig.~\ref{fig:crash_types}). SCP+MPPI has 14 unwarned violations in 880
episodes, against 6 and 7 for the 512-sequence samplers and 7 and 13 for the time-matched ones. That stronger
optimisation, by SCP or by more samples, leads to more violations the model did not foresee would be consistent with
it planning closer to where the model is wrong, but paired by seed none of these differences is significant
($p\ge0.12$).

% ---------------------------------------------------------------------------------------------------------------------
\subsection{When violations happen}\label{app:survival}

\begin{figure}[htbp]
  \centering
  \includegraphics[width=\linewidth]{\figdir/appendix/fig_survival_pg1_c1_5.pdf}
  \caption{Share of episodes without a violation by each step, at each margin, on PointGoal1's $c=1.5$ model at
  $N=10$. Episodes end early only at a violation, so these are Kaplan--Meier curves without censoring.}
  \label{fig:survival}
\end{figure}

Most violations happen within the first 300 or so steps, after which the curves flatten (Fig.~\ref{fig:survival});
pure SCP's keeps falling for longer.
 in the appendix. Files and numbers as in experiments_results.tex.
% Order is by importance; drop subsections from the end if the appendix must be short.
% =====================================================================================================================
\providecommand{\figdir}{../paper_graphs}

\section{Additional Results}\label{app:results}

% ---------------------------------------------------------------------------------------------------------------------
\subsection{Per-model results}\label{app:per_model}

\begin{table}[htbp]
  \centering
  \caption{Table~\ref{tab:pooled} on the $c=1.2$ world models.}
  \label{tab:pooled_c1_2}
  \small
  \input{\figdir/appendix/table_pooled_c1_2}
\end{table}

\begin{table}[htbp]
  \centering
  \caption{Table~\ref{tab:pooled} on the unbounded ($c=\infty$) world models.}
  \label{tab:pooled_cinf}
  \small
  \input{\figdir/appendix/table_pooled_cinf}
\end{table}

Table~\ref{tab:per_model} gives each planner's violation rate and return on every $N=10$ world model, pooled over
$\epsilon\ge0.1$, and Figs.~\ref{fig:grid_violations} and~\ref{fig:grid_return} show the full sweeps. Pure SCP is
significantly less safe than SCP+MPPI on every model. Against the samplers, the only significant difference is MPPI
with 512 sequences on PointGoal2 with $c=1.2$, which violates less.

\begin{table}[htbp]
  \centering
  \caption{Each planner on each $N=10$ world model, pooled over $\epsilon\ge0.1$: violation rate (Wilson) and return
  (bootstrap) with 95\% intervals, and McNemar's $p$ for SCP+MPPI against it on paired seeds, Holm-corrected over the
  six models.}
  \label{tab:per_model}
  \small
  \input{\figdir/appendix/table_per_model}
\end{table}

\begin{figure}[htbp]
  \centering
  \includegraphics[width=\linewidth]{\figdir/appendix/fig_eps_sweep_all_bounds_violations.pdf}
  \caption{Violation rate against $\epsilon$ for every planner on every $N=10$ world model (rows: tasks; columns:
  Lipschitz bound $c$), with 95\% Wilson intervals.}
  \label{fig:grid_violations}
\end{figure}

\begin{figure}[htbp]
  \centering
  \includegraphics[width=\linewidth]{\figdir/appendix/fig_eps_sweep_all_bounds_return.pdf}
  \caption{Return against $\epsilon$ for every planner on every $N=10$ world model, with 95\% bootstrap intervals.}
  \label{fig:grid_return}
\end{figure}

% ---------------------------------------------------------------------------------------------------------------------
\subsection{Lipschitz bound and SCP's convergence}\label{app:lipschitz}

\begin{figure}[htbp]
  \centering
  \includegraphics[width=\linewidth]{\figdir/appendix/fig_lipschitz_convergence.pdf}
  \caption{Share of SCP+MPPI's steps that reached the 40-iteration cap against $\epsilon$, for each Lipschitz bound,
  at $N=10$.}
  \label{fig:lipschitz_convergence}
\end{figure}

Without a Lipschitz bound, SCP+MPPI reaches the iteration cap on more steps than with either bound at every margin,
except at $\epsilon=0.5$ on PointGoal2, where $c=1.5$ and $c=\infty$ are within 0.02 (Fig.~\ref{fig:lipschitz_convergence}).
Which bound converges best differs by task: $c=1.5$ on PointGoal1, $c=1.2$ on PointGoal2.

% ---------------------------------------------------------------------------------------------------------------------
\subsection{What more samples buy at \texorpdfstring{$N=10$}{N=10}}\label{app:samples}

\begin{figure}[htbp]
  \centering
  \includegraphics[width=\linewidth]{\figdir/ablation/fig_sample_vs_time_matched.pdf}
  \caption{Fig.~\ref{fig:eps_sweep} with the samplers at both budgets: 512 sequences (open markers, dashed) and
  time-matched (filled, solid).}
  \label{fig:sample_vs_time_matched}
\end{figure}

\begin{figure}[htbp]
  \centering
  \includegraphics[width=\linewidth]{\figdir/ablation/fig_return_vs_violations.pdf}
  \caption{Return against violation rate on each task's $c=1.5$ model at $N=10$, pooled over $\epsilon\ge0.1$, with
  95\% intervals. Arrows lead from each sampler with 512 sequences to the same sampler time-matched.}
  \label{fig:trade_off}
\end{figure}

Figs.~\ref{fig:sample_vs_time_matched} and~\ref{fig:trade_off} show what 128 times as many sequences buy the samplers.
On PointGoal1's $c=1.5$ model, the time-matched samplers earn more than with 512 sequences but violate more often
(MPPI: return 0.78 to 1.15, violation rate 0.19 to 0.24; CEM: 0.84 to 0.94 and 0.14 to 0.26), while SCP+MPPI reaches
the higher return at the lower violation rate (1.15 and 0.17). On PointGoal2 the extra sequences raise return (0.35 to
0.51 and 0.49) without changing the violation rate. These shifts lie within the intervals; the paired tests of
Sec.~\ref{sec:q1} find no significant difference in safety between SCP+MPPI and either budget.

% ---------------------------------------------------------------------------------------------------------------------
\subsection{All budgets at every horizon}\label{app:horizon_budgets}

\begin{figure}[htbp]
  \centering
  \includegraphics[width=\linewidth]{\figdir/ablation/fig_horizon_all_budgets.pdf}
  \caption{Fig.~\ref{fig:horizon} with the samplers at both 512 sequences (open markers, dashed) and their
  time-matched budget (filled, solid).}
  \label{fig:horizon_all}
\end{figure}

With 512 sequences the samplers' return at $N=20$ is close to SCP+MPPI's (MPPI 2.07, CEM 2.58, against 2.34; neither
difference significant), and at $N=30$ below it (1.59 and 1.22 against 2.26; $p=0.043$ and $0.013$)
(Fig.~\ref{fig:horizon_all}). Together with Fig.~\ref{fig:compute}, this shows that the samplers' best budget depends on
the sampler and the horizon.

% ---------------------------------------------------------------------------------------------------------------------
\subsection{Does optimisation exploit the model?}\label{app:optimism}

\begin{figure}[htbp]
  \centering
  \includegraphics[width=\linewidth]{\figdir/ablation/fig_model_optimism.pdf}
  \caption{Model optimism: each executed plan's predicted return (discounted predicted rewards plus terminal value)
  minus the discounted return received from that step on, averaged over each episode, with 95\% bootstrap intervals.
  $N=10$ pools every world model ($\epsilon\ge0.1$); $N=20$ and $30$ are the horizon study's models.}
  \label{fig:optimism}
\end{figure}

A gradient-based optimiser might exploit model error, finding plans whose predicted return the environment does not
deliver. We measure each executed plan's optimism: its predicted return, whose terminal value carries the goal bonus,
minus the discounted return actually received from that step on, using steps at least 299 steps before the episode
ends so that the part of the future the episode did not reach weighs under 5\%. Every planner is optimistic by about
2, as expected of a value fitted to the 0.7 expectile, but by nearly the same amount (Fig.~\ref{fig:optimism}): 2.06
(pure SCP) to 2.19 at $N=10$, 2.20 to 2.44 at $N=20$ and 2.41 to 2.56 at $N=30$, with SCP+MPPI inside the samplers'
range at every horizon and no trend with the sampling budget. We find no evidence that SCP, or a larger sampling
budget, exploits the model more than sampling with 512 sequences does.

% ---------------------------------------------------------------------------------------------------------------------
\subsection{Warned and unwarned violations}\label{app:crashes}

\begin{figure}[htbp]
  \centering
  \includegraphics[width=\linewidth]{\figdir/ablation/fig_crash_types.pdf}
  \caption{Violations per 100 episodes over the six $N=10$ models and $\epsilon\ge0.1$, split into those preceded
  within $N$ steps by a plan that broke the barrier in the model (warned, light) and those that were not (unwarned,
  dark, counts above the bars).}
  \label{fig:crash_types}
\end{figure}

Few violations come without warning (Fig.~\ref{fig:crash_types}). SCP+MPPI has 14 unwarned violations in 880
episodes, against 6 and 7 for the 512-sequence samplers and 7 and 13 for the time-matched ones. That stronger
optimisation, by SCP or by more samples, leads to more violations the model did not foresee would be consistent with
it planning closer to where the model is wrong, but paired by seed none of these differences is significant
($p\ge0.12$).

% ---------------------------------------------------------------------------------------------------------------------
\subsection{When violations happen}\label{app:survival}

\begin{figure}[htbp]
  \centering
  \includegraphics[width=\linewidth]{\figdir/appendix/fig_survival_pg1_c1_5.pdf}
  \caption{Share of episodes without a violation by each step, at each margin, on PointGoal1's $c=1.5$ model at
  $N=10$. Episodes end early only at a violation, so these are Kaplan--Meier curves without censoring.}
  \label{fig:survival}
\end{figure}

Most violations happen within the first 300 or so steps, after which the curves flatten (Fig.~\ref{fig:survival});
pure SCP's keeps falling for longer.
